% !TeX root = ../informe.tex
% ===========================================================================
%  Portada. Los datos se toman de config/datos.tex — no editar aquí.
% ===========================================================================

% Imprime un integrante (nombre + correo). Si el nombre está vacío, no
% imprime nada: así sobran entradas en datos.tex sin dejar huecos.
\newcommand{\integrante}[2]{%
    \ifx\relax#1\relax\else
        {\large #1\par}%
        {\small\ttfamily #2\par}%
        \vspace{8pt}%
    \fi
}

\begin{titlepage}
\centering
\thispagestyle{empty}
% Interlineado sencillo solo en la portada: con el 1.5 del cuerpo, cuatro
% integrantes con correo desbordan a una segunda página.
\singlespacing

% Descomentar cuando se tenga el logotipo en figuras/logo.png
% \includegraphics[width=0.28\textwidth]{logo.png}\\[1.5cm]

{\scshape\large \nombreUniversidad\par}
\vspace{0.4cm}
{\scshape\large \nombrePrograma\par}
\vspace{0.2cm}
{\large \nombreCurso\par}

\vspace{1.6cm}

\rule{\textwidth}{0.4pt}\\[0.5cm]
{\huge\bfseries \tituloInforme\par}
\vspace{0.4cm}
{\Large \subtituloInforme\par}
\rule{\textwidth}{0.4pt}\\[0.5cm]

{\large \tipoTrabajo\par}

\vfill

{\large\bfseries Autores\par}
\vspace{0.4cm}

\integrante{\autorUno}{\correoUno}
\integrante{\autorDos}{\correoDos}
\integrante{\autorTres}{\correoTres}
\integrante{\autorCuatro}{\correoCuatro}

\vspace{0.4cm}

{\large\bfseries Docente\par}
\vspace{0.25cm}
{\large \nombreDocente\par}

\vfill

{\large \ciudadFecha\par}

\end{titlepage}
